% ---
% id: reductions/rem:reading-master
% kind: rem
% title: "The operator form of the branch-weight constraint"
% summary: >-
%   The branch weights of [[lem:master]] (iv) are the only constraint the
%   branching rule produces, and the operator form is the one to impose: it
%   fixes each block's weight conditioned on everything Alice and Eve hold,
%   while its trace alone admits a block index perfectly correlated with Eve.
% prominence: standard
% status: accepted
% uses: [cor:transported, lem:master]
% assumes: []
% refs: []
% provenance: newly-derived
% lean:
% ---

\begin{remark}
\label{rem:reading-master}
Part~(iv) of Lemma~\ref{lem:master} is the only place the branching rule
produces a constraint, and the operator form is the one worth imposing: it
says the weight of block $j$ is $q_j$ \emph{conditioned on everything Alice
and Eve hold}, and Corollary~\ref{cor:transported}\,(i) turns it into a sound
constraint, one per round. The scalar $\tr[\Pi_j\sigma]=q_j$ is only its trace
and is strictly weaker: it still admits a squashed state whose block index is
perfectly correlated with Eve's register. Nothing is averaged anywhere in the
rule, because averaging over $j$ before squashing would lose the branch
structure and with it the relation.
\end{remark}
